\documentclass[10pt]{article}
\usepackage[margin=0.8in]{geometry}
\usepackage{amsmath,amssymb,amsthm,mathtools,booktabs,microtype,hyperref}
\hypersetup{hidelinks,pdftitle={A canonical symmetry excludes the ramified quartic completion}}
\newtheorem{theorem}{Theorem}
\newtheorem{lemma}[theorem]{Lemma}
\newcommand{\R}{\mathbb R}
\newcommand{\Cr}{\mathcal C_{\mathrm{ram}}}
\newcommand{\Cc}{\mathcal C_{\mathrm{can}}}
\newcommand{\tr}{\operatorname{tr}}
\newcommand{\ip}[2]{\langle #1,#2\rangle}
\title{A canonical symmetry excludes the ramified quartic completion}
\author{}\date{}
\begin{document}
\maketitle\vspace{-2.5em}
\begin{abstract}
An explicit all-positive valid source has no matrix selector polynomial
in the square root of the Bob half-angle coordinate. Four positive-definite integer block
Hankel matrices give an exact separating inequality. A physical exchange
of Alice and Bob transports a source inside the ramified completion cone
to this excluded source. The same transformation moves Bob quarter-power
testing functions to Alice and produces two incomparable certificate cones.
\end{abstract}

\section{The counterexample and its validity}

Use the source order
$c=(d,\lambda,n,m,C,A,D,\alpha,\beta)\in\R_+^9$ and the complete
canonical parent $F_c$. The cone $\Cc$ consists of sources valid at every
state and every pair of projective measurement angles. The
two-completion matrices are
\[
 G_J=\begin{pmatrix}P+K&J\\J^T&N+K\end{pmatrix},\qquad
 H_J=\begin{pmatrix}P+E&K-E-J\\K-E-J^T&N+E\end{pmatrix}.
\]
Here $P,N,E,K$ are exactly the source-affine quadratic matrices in the
canonical completion note~\cite{accepted}. The fixed cone $\Cr$ requires
$v=16s^2$, $s\ge0$, and one matrix polynomial
$J(s)=\sum_{k=0}^4J_ks^k$ making both matrices positive semidefinite.
The resulting Gram pencil has order 49.

The explicit Gram matrices in Ref.~\cite{accepted} give
\[
 \bar c=\left(1,2^{-34},2^{-14},2^{-24},\frac1{576},161,
                   \frac12,76,\frac1{250}\right)\in\Cr\subseteq\Cc.
 \tag{1}
\]

\begin{theorem}\label{thm:main}
The all-positive rational source
\[
 \boxed{c^\dagger=\left(1,\frac1{576},161,2^{-24},2^{-34},2^{-14},
                   \frac12,\frac1{250},76\right)}           \tag{2}
\]
belongs to $\Cc\setminus\Cr$. Consequently
$\Cr\subsetneq\Cc$. The source $c^\dagger$ admits no selector polynomial
in $s=\sqrt v/4$, of any degree. The order-392 eight-copy outer lift is
strictly larger than $Q_{222}$ by the polar assembly theorem~\cite{assembly}.
\end{theorem}

Define the involution
\[
 \pi c=(d,C,A,m,\lambda,n,D,\beta,\alpha),\qquad
 S(W,X,Y,Z)=(W,Y,X,Z).
\]
The full parent satisfies the exact identity
\[
 \boxed{u^2v^2F_{\pi c}(-1/v,-1/u;S\psi)=F_c(u,v;\psi).} \tag{3}
\]
The eight event packets, in the source
order $(\lambda,n,m,C,A,D,\alpha,\beta)$, are carried respectively to
$(C,A,m,\lambda,n,D,\beta,\alpha)$. The $m,D$ roots change sign.
The $d$ coefficient is unchanged: its three positive packets are permuted
or fixed, and $D_uD_vZ^2$ is unchanged after clearing the angles.
Thus (3) holds for every source.

With homogeneous angle pairs
$a=(a_0,a_1)$, $b=(b_0,b_1)$, identity (3) is the polynomial identity
\[
 \boxed{F^h_{\pi c}\bigl((b_1,-b_0),(a_1,-a_0);S\psi\bigr)
                       =F^h_c(a,b;\psi).}                 \tag{4}
\]
It holds on the complete product of angle circles, including literal zero
coordinates. It permutes zero source weights exactly, keeps all state
signs, and acts separately on source and testing coordinates. Real and
imaginary state parts give the same conclusion for complex states.
Thus $\pi\Cc=\Cc$. Since $c^\dagger=\pi\bar c$, (1) proves validity
of (2). It remains to exclude its quartic completion.

\section{An exact separating trace identity}

For $\sigma=G,H$, let $N_{\sigma,0},\ldots,N_{\sigma,4}$ be the ten
\emph{integer} symmetric $4\times4$ matrices in Appendix~\ref{app:data}.
Form the four integer block Hankel matrices
\[
 \mathsf U_\sigma=(N_{\sigma,i+j})_{i,j=0}^{2},\qquad
 \mathsf V_\sigma=(N_{\sigma,i+j+1})_{i,j=0}^{1}.           \tag{5}
\]
They have orders $12,8,12,8$ and are positive definite. Moreover,
\[
 (N_{G,k})_{\{1,2\},\{3,4\}}
 =(N_{H,k})_{\{1,2\},\{3,4\}}
 \qquad(0\le k\le4).                                    \tag{6}
\]
Positive definiteness follows from the rational pivots in
Appendix~\ref{app:data}; the upper-right blocks agree entry by entry.

Write $Q(s)=\sum_{k=0}^4Q_ks^k$ for either completed matrix. Its
half-line Gram representation is
\[
 Q(s)=Z_2(s)^TA_QZ_2(s)+sZ_1(s)^TB_QZ_1(s),\quad
 Z_2=(I_4,sI_4,s^2I_4)^T,\quad Z_1=(I_4,sI_4)^T,
\]
where $A_Q,B_Q\succeq0$. Expanding gives
\[
 \sum_{k=0}^4\ip{N_{\sigma,k}}{Q_k}
 =\ip{\mathsf U_\sigma}{A_Q}+\ip{\mathsf V_\sigma}{B_Q}.
                                                               \tag{7}
\]
Each $J_k$ occurs in opposite off-diagonal blocks of $G_k$ and $H_k$,
so (6) cancels every selector coefficient. Consequently the sum of (7)
is source-affine and nonnegative for every feasible certificate, including
zero weights and singular Grams.

The resulting integer functional in the stated source order is
\[
\begin{array}{c|r@{\qquad}c|r@{\qquad}c|r}
 d&-97633735551856&\lambda&24154193783510329&n&85793705264\\
 m&1787034292680750290&C&3303160925893294&A&16956553777510394\\
 D&13682527570574&\alpha&670221724556174&\beta&366431361010
\end{array}
                                                               \tag{8}
\]
Denote the functional with these coefficients by $\ell(c)$.
Substituting $P,N,E,K$, with $v=16s^2$, gives
the exact identity
\[
 \boxed{12\ell(c)=
 \ip{\mathsf U_G}{A_G}+\ip{\mathsf V_G}{B_G}
 +\ip{\mathsf U_H}{A_H}+\ip{\mathsf V_H}{B_H}.}             \tag{9}
\]
Only the moments of degrees $0,2,4$ contribute to the nine source
coefficients. If $Q_j^{(v)}(c,0)$ is the coefficient of $v^j$ with $J=0$,
the left side of (9) is
\[
 \sum_{j=0}^2 16^j\bigl[
 \ip{N_{G,2j}}{G_j^{(v)}(c,0)}+
 \ip{N_{H,2j}}{H_j^{(v)}(c,0)}\bigr].                     \tag{10}
\]
Thus $\ell\ge0$ on $\Cr$. But direct rational evaluation yields
\[
 \boxed{\ell(c^\dagger)=
 -\frac{32607070960447356138172589}{9663676416000}<0.}      \tag{11}
\]
The four PSD trace pairings in (9) cannot have a negative sum. This
excludes the quartic selector.

To exclude every polynomial degree in the same coordinate, suppose that
a polynomial matrix $J(s)$ made both completions PSD on $s\geq0$.
Their principal minors would give
\[
 |J_{ij}(s)|^2\leq(P+K)_{ii}(16s^2)(N+K)_{jj}(16s^2)=O(s^8).
\]
The source diagonal blocks have degree at most four in $s$, so
$\deg J_{ij}\leq4$. Every quartic matrix polynomial PSD on the half-line
has the displayed Gram representation: factor the even globally PSD
matrix polynomial $Q(t^2)$, average its polynomial SOS at $t$ and $-t$,
and separate the even and odd roots. With $s=t^2$, the even roots have the form
$B_0+sB_1+s^2B_2$ and the odd roots have the form
$t(C_0+sC_1)$, yielding $Z_2^TA_QZ_2+sZ_1^TB_QZ_1$.
Thus every polynomial selector would give the same four Gram blocks,
already excluded by (11). This proves Theorem~\ref{thm:main}.

\section{The contact mechanism and the remaining problem}

The separation already holds at the exact contact source $\pi c_*$
underlying (1). Its contact in positive angle gauges is
\[
 u=\frac1{16},\qquad v=13,\qquad
 \psi=\left(140,-\frac1{50},\frac1{32},1\right),
\]
and
\[
 \ell(\pi c_*)=
 -\frac{4844338143654932884745484309348314410}
        {802811514300974151343419}<0.                    \tag{12}
\]
In the full state variables, write
\[
 \eta=L(1/16)\psi=
 \left(\frac{71681}{8192},\frac{17}{6400},
       \frac{279}{32},-\frac{801}{800}\right).
\]
At $U=1/256$ and $s=\sqrt{13}/4$, any putative completion must satisfy
\[
 G(s)(\eta_1,\eta_2,U\eta_3,U\eta_4)^T=0,\qquad
 H(s)\eta^T=0.                                          \tag{13}
\]
These follow from the split identity and positivity of both
summands at a zero. The separator excludes such a global quartic
completion and remains negative after opening both zero event weights to
the all-positive source (2).

The involution $\pi$ preserves canonical validity and transports its
testing functions between the parties. The five compact Bob weights in
affine coordinates are
\[
 \frac{(1,\sqrt{|v|}/4,|v|/16,|v|^{1/4}/2,|v|^{3/4}/8)}
      {\sqrt{1+v^2}}.
\]
Under (3), they become the fixed Alice weights
\[
 \boxed{
 \frac{(|u|,\sqrt{|u|}/4,1/16,|u|^{3/4}/2,|u|^{1/4}/8)}
      {\sqrt{1+u^2}}.}                                  \tag{14}
\]
Their branches and normalization depend only on the test. The homogeneous
pullback of the forty functions represents $b_0^2F^h_{\pi c}$,
where the original family represented $a_1^2F^h_c$; both multipliers are
nonnegative and positive on a dense testing chart.

The transported functions lie outside the original testing span. Restrict to
$u>0$, $v=1$, $\psi=(1,0,0,0)$. Every old compact feature is a constant
times a polynomial of degree at most three divided by $(1+u^2)^{3/2}$,
and is analytic at $u=0$. One transported feature is
\[
 -\frac{u^{1/4}}{16\sqrt2\sqrt{1+u^2}}.
\]
It has nonintegral order of vanishing and cannot lie in that span.
Over the common multiplier $a_1^2b_0^2$, multiply old roots by $b_0$ and
transported roots by $a_1$; the same feature becomes
$-u^{5/4}/(16\sqrt2(1+u^2))$, so the distinction persists.

In particular $\Cr$ and $\pi\Cr$ are incomparable: (1), (2), and
$\pi^2=1$ give a point in each set difference. It remains open whether
$\Cr+\pi\Cr=\Cc$, or whether further testing functions are required.

A finite construction of the full cone requires a source-independent
algebraic testing space whose Gram cone contains every valid slack,
possibly after one fixed pure-testing multiplier. The contact and
conductor identities and the qualitative tensor criterion of
Ref.~\cite{companion} do not specify such a global space.

\begin{thebibliography}{3}\small
\bibitem{accepted} \emph{Canonical Bell validity through two matrix inequalities};
\emph{An exact obstruction to the quadratic half-line completion},
and the ramified order-49 data.
\href{https://github.com/anubhav70747/q222-finite-sdp/tree/16b37f6deb198fa48134145b4d12b7cc7d37be36/extracted-lifts/canonical-completion}{Public repository, canonical-completion, commit 16b37f6}.
\bibitem{assembly} \emph{Q222 Polar Assembly},
\href{https://github.com/anubhav70747/q222-finite-sdp/tree/16b37f6deb198fa48134145b4d12b7cc7d37be36/extracted-lifts/canonical-completion/full-polar}{same commit, full-polar}.
\bibitem{companion} Reconstructed R21 companion,
\href{https://github.com/anubhav70747/q222-finite-sdp/tree/16b37f6deb198fa48134145b4d12b7cc7d37be36/latex}{same commit, latex},
Chapters 3--6; in particular TF12, SG2, SG7, and CC1--CC6.
\end{thebibliography}

\clearpage
\appendix
\section{The complete integer moment certificate}\label{app:data}
All matrices below are used as integers in (5)--(10).
The normalized moment matrices are $N/(15\cdot2^{60})$; equations
(5)--(10) use their integer numerators.

For each matrix $R$ in (5), set
\[
 p_j=R_{jj}-\sum_{k<j}L_{jk}^2p_k,\qquad
 L_{ij}=\frac{R_{ij}-\sum_{k<j}L_{ik}p_kL_{jk}}{p_j},
 \quad L_{jj}=1\quad(i>j).
\]
It gives $R=L\operatorname{diag}(p_j)L^T$. The exact floors of the
smallest pivots are
\[
\begin{array}{c|rrrr}
 R&\mathsf U_G&\mathsf V_G&\mathsf U_H&\mathsf V_H\\\hline
 \left\lfloor\min_jp_j\right\rfloor&794&12397&67408&1527738.
\end{array}
\]
These positive pivot bounds prove that all four matrices are positive definite. The shared
upper-right blocks in each pair
$N_{G,k},N_{H,k}$ make (6) directly visible.

\[
N_{G,0}=\begin{pmatrix}
72283584904157745 & 98314305759870 & 214652490042795 & -146888491666410\\
98314305759870 & 135542238435 & 268778302935 & -201769656060\\
214652490042795 & 268778302935 & 1297549019460 & -401316089475\\
-146888491666410 & -201769656060 & -401316089475 & 301979389605
\end{pmatrix}.
\]

\[
N_{G,1}=\begin{pmatrix}
29553765112790565 & 39441163448385 & 105424012102020 & -58918574211480\\
39441163448385 & 54072935055 & 110357620980 & -80768779620\\
105424012102020 & 110357620980 & 1079281903875 & -164509196205\\
-58918574211480 & -80768779620 & -164509196205 & 120695415960
\end{pmatrix}.
\]

\[
N_{G,2}=\begin{pmatrix}
12789595044635020 & 15916490027610 & 66271058811705 & -23763057906660\\
15916490027610 & 21653142535 & 47712869220 & -32349910290\\
66271058811705 & 47712869220 & 1287764254705 & -70816288035\\
-23763057906660 & -32349910290 & -70816288035 & 48333328645
\end{pmatrix}.
\]

\[
N_{G,3}=\begin{pmatrix}
6842765115107160 & 6582732635010 & 46453400242755 & -9802159230735\\
6582732635010 & 8693049165 & 22782075435 & -12984692925\\
46453400242755 & 22782075435 & 1981266595425 & -33639345030\\
-9802159230735 & -12984692925 & -33639345030 & 19395668925
\end{pmatrix}.
\]

\[
N_{G,4}=\begin{pmatrix}
20910075570761184 & 4127281354275 & 33265473661770 & -6066895933245\\
4127281354275 & 3608024196 & 12294756105 & -5380775115\\
33265473661770 & 12294756105 & 5215970784456 & -19334942190\\
-6066895933245 & -5380775115 & -19334942190 & 8026177596
\end{pmatrix}.
\]

\clearpage
\section*{The five $H$ moments}
\[
N_{H,0}=\begin{pmatrix}
51123677524583550 & -5656910689200 & 214652490042795 & -146888491666410\\
-5656910689200 & 10186908915 & 268778302935 & -201769656060\\
214652490042795 & 268778302935 & 146089272033990 & -170629846948365\\
-146888491666410 & -201769656060 & -170629846948365 & 2302846743948255
\end{pmatrix}.
\]

\[
N_{H,1}=\begin{pmatrix}
765887260600020 & 39305335635 & 105424012102020 & -58918574211480\\
39305335635 & 2703893715 & 110357620980 & -80768779620\\
105424012102020 & 110357620980 & 92030582539560 & -33880347103335\\
-58918574211480 & -80768779620 & -33880347103335 & 160298762935095
\end{pmatrix}.
\]

\[
N_{H,2}=\begin{pmatrix}
89446398381610 & 50022371745 & 66271058811705 & -23763057906660\\
50022371745 & 931545700 & 47712869220 & -32349910290\\
66271058811705 & 47712869220 & 64202652830605 & -13709279518560\\
-23763057906660 & -32349910290 & -13709279518560 & 18313581996550
\end{pmatrix}.
\]

\[
N_{H,3}=\begin{pmatrix}
48316110408750 & 24625867005 & 46453400242755 & -9802159230735\\
24625867005 & 328684920 & 22782075435 & -12984692925\\
46453400242755 & 22782075435 & 45819719138655 & -8466528163695\\
-9802159230735 & -12984692925 & -8466528163695 & 3366601507110
\end{pmatrix}.
\]

\[
N_{H,4}=\begin{pmatrix}
33727791386796 & 13046747130 & 33265473661770 & -6066895933245\\
13046747130 & 116838834 & 12294756105 & -5380775115\\
33265473661770 & 12294756105 & 32917216546116 & -5858691614355\\
-6066895933245 & -5380775115 & -5858691614355 & 1299098148456
\end{pmatrix}.
\]

\end{document}
